\section{Introduction}\label{sec:introduction}

A random regular graph is sparse even when its order tends to infinity:
each row of its adjacency matrix has exactly \(d\) nonzero entries.
Nevertheless, its bulk eigenvalues are expected to exhibit the local
statistics of a real symmetric Gaussian matrix.
We prove this prediction for every fixed degree \(d\ge3\), at each fixed
energy in the interior of the limiting spectrum.

\subsection{Main result}

For an admissible integer \(n\), let \(G_{n,d}\) be uniform among the
simple labelled \(d\)-regular graphs on \(\{1,\ldots,n\}\).
Here admissible means that this set is nonempty.
Let \(A_{n,d}\) be its unweighted adjacency matrix, with eigenvalues
\(\lambda_1,\ldots,\lambda_n\), counted with multiplicity.
No conditioning on connectedness, bipartiteness, or another graph
property is imposed.

The limiting density is the Kesten--McKay density
\[
 \rho_d(t)=\frac{d\sqrt{4(d-1)-t^2}}{2\pi(d^2-t^2)}
 \qquad \bigl(|t|<2\sqrt{d-1}\bigr).
\]
For a fixed energy \(E\) in this interval, define the unfolded
point measure
\[
 \Xi_{n,d,E}=\sum_{i=1}^n
       \delta_{\,n\rho_d(E)(\lambda_i-E)}.
\]
For comparison, let \(H_n\) be a normalized GOE matrix: its entries on
and above the diagonal are independent centered Gaussians, with
variance \(1/n\) off the diagonal and \(2/n\) on the diagonal.
If \(\mu_1,\ldots,\mu_n\) are its eigenvalues, put
\[
 \Xi_n^{\mathrm{GOE}}
       =\sum_{i=1}^n\delta_{(n/\pi)\mu_i}.
\]

\newpage
\begin{theorem}\label{thm:main}
For every fixed integer \(d\ge3\), every fixed
\(E\in(-2\sqrt{d-1},2\sqrt{d-1})\), and every nonnegative continuous
compactly supported \(h:\R\to\R\),
\[
 \left|
 \EE\exp\!\left(-\int h\dd\Xi_{n,d,E}\right)
 -
 \EE\exp\!\left(-\int h\dd\Xi_n^{\mathrm{GOE}}\right)
 \right|\longrightarrow0
\]
as \(n\to\infty\) through all admissible sizes.
\end{theorem}

The GOE reference point measures converge in law for the vague
topology \cite{ValkoVirag}.
Thus Theorem~\ref{thm:main} identifies the full local point-process
limit, not only its one-point density or an energy-averaged statistic.
It proves the fixed-degree bulk-universality conjecture stated in
\cite[Conjecture 2.11]{BourgadeHuang}, in a fixed-energy
Laplace-functional formulation.

\subsection{Background and significance}

The density above originates in the spectral measure of the regular
tree and the limiting spectrum of random regular graphs
\cite{Kesten,McKay}.
It describes eigenvalue counts on macroscopic intervals, but does not
determine the spacings between consecutive eigenvalues.
Jakobson, Miller, Rivin, and Rudnick
\cite[Sections~1 and~5]{JakobsonEtAl} formulated the GOE prediction
for fixed-degree regular graphs and gave numerical evidence for it.
Their empirical nearest-neighbor spacing formulation precedes the
fixed-energy point-process formulation studied here.
The prediction is particularly striking at fixed degree: local tree
geometry persists, whereas the proposed eigenvalue statistics agree
with those of a dense Gaussian matrix.

The development of local laws separated spectral control from the
identification of microscopic statistics. Bauerschmidt, Knowles, and
Yau \cite{BauerschmidtKnowlesYauLocal} established a local semicircle
law for growing degrees using local resampling. Bauerschmidt, Huang,
Knowles, and Yau \cite[Theorems 1.1 and 1.2]{BauerschmidtEtAl}
proved bulk gap universality and energy-averaged correlation
universality in a polynomial degree range. More recently, Bourgade
and Huang \cite[Proposition 2.10]{BourgadeHuang} obtained fixed-energy
point-process universality for
\((\log n)^{24}\ll d\le n^{1/2}\) by loop equations, and stated the
remaining fixed-degree problem in their Conjecture 2.11.

For fixed degree the local reference law is Kesten--McKay.
Bauerschmidt, Huang, and Yau \cite{BauerschmidtHuangYauLocal}
established the local law and delocalization for sufficiently large
fixed degrees. Huang and Yau \cite{HuangYau} subsequently proved
rigidity, complete delocalization, and local resolvent estimates for
every fixed \(d\ge3\). Their estimates supply the spectral input for
this paper. They control smoothed eigenvalue counts and typical tree
neighborhoods, while the conclusion sought here concerns unsmoothed
statistics at one fixed energy on the scale of an individual spacing.

For every fixed \(d\ge3\), Huang, McKenzie, and Yau
\cite[Theorems~1.1 and~1.2]{HuangMcKenzieYau}
proved optimal eigenvalue rigidity, up to subpolynomial factors,
and GOE edge universality.
Their edge result identifies the joint fluctuations of any fixed
number of extreme nontrivial eigenvalues at either spectral edge.

Eigenvector Gaussianity provides the link we need from spectral control
to microscopic bulk statistics.
The Gaussian-wave model on a regular tree appears in
\cite{ElonEigenvectors,CsokaEtAl}.
Backhausz and Szegedy \cite[Theorems 2.1 and 2.2]{BackhauszSzegedy}
proved qualitative Gaussian-wave limits for almost eigenvectors of
random regular graphs, allowing a variance factor between zero and
one in the weak limit. Their argument uses entropy inequalities for
local observations; the underlying colored-star counting framework
was developed in \cite{BackhauszSzegedyTrees} and further studied in
\cite{BackhauszBordenaveSzegedy}. Recent developments include
variance-one Gaussian waves for spectral-edge eigenvectors
\cite{HeHuangYauWaves} and an extension of the entropy framework to
graphs of finite cone type \cite{DemboMcKenzie}.
Our argument develops this entropy method into a different input:
polynomially accurate, variance-one joint moments at sampled graph
edges. The estimates hold conditional on each graph in an event of
probability tending to one, uniformly over finitely many chosen bulk
eigenvectors. Weak convergence alone does not provide
these high-moment estimates. The proof below supplies them through
a strict tensor estimate and a quantitative Hermite expansion.

Reversible edge switches are also established tools in local laws
and universality; see \cite[Section 3]{BauerschmidtEtAl} and
\cite[Section 7.2]{BauerschmidtHuangYauLocal}.
We use Gaussian endpoint marks to extend the discrete symmetry to
an orthogonal action on an exact microscopic limiting law. Increasing
the number of sampled edge pairs within that limit then inserts a Gaussian
matrix perturbation. The last step uses the local relaxation of
Dyson's matrix evolution \cite{DysonBrownian,LandonYau}, in the
fixed-energy form proved by Landon, Sosoe, and Yau
\cite{LandonSosoeYau}. To obtain the full point-process conclusion
from fixed-order correlations, we also give a compact-count bound
using the Gaussian superposition and decimation identity of
Forrester and Rains \cite{ForresterRains}.

\subsection{Proof strategy}

We first work in the coordinates \(x_i=n(\lambda_i-E)\), before the
last density factor \(\rho_d(E)\) is applied.
For a real eigenvector normalized to have squared norm \(n\), sampling
its values at the two ends of a uniform oriented edge produces a pair
of random real numbers. We call each sampled edge and its two endpoint
coordinates a channel.
The first part of the proof shows that, for any fixed number of
eigenvectors near \(E\), every fixed joint moment of these sampled
values is polynomially close to the corresponding Gaussian moment,
with probability tending to one over the graph.
The comparison consists of centered, unit-variance Gaussian pairs,
independent across the eigenvectors, with covariance \(E/d\) within each pair.

The following steps then identify every subsequential microscopic
point-process limit.

\begin{enumerate}
\item \emph{Quantitative Gaussian marks.}
Section~\ref{sec:entropy} proves the endpoint moment statement.
A colored-star counting inequality becomes an entropy inequality after
the eigenvectors are smoothed by a Gaussian field.
The tree logarithmic potential evaluates the entropy cost.
A strict tensor inequality turns the resulting entropy deficit into
Gaussian moment matching by induction over Hermite degree.

\item \emph{Microscopic limits with controlled tails.}
Section~\ref{sec:limits} uses a two-edge switch and the moment statement
to control resolvents at large fixed microscopic heights.
It then extracts the point process together with its Gaussian endpoint
marks. The limiting resolvent has a symmetric principal-value
representation, and the point counts have a power-saving discrepancy
from constant density.

\item \emph{Continuous switching and Gaussian insertion.}
Section~\ref{sec:switching} combines the discrete switch symmetry with
orthogonal invariance of the Gaussian marks.
This produces exact continuous switching identities for the
microscopic limiting law.
After two controlled truncations, many weak switches give finite
diagonal matrices with additive GOE noise whose eigenvalue point
processes converge to that same limiting law.

\item \emph{Identification at one energy.}
Section~\ref{sec:dbm} verifies the hypotheses of fixed-energy Dyson
Brownian motion universality for these diagonal matrices.
An exact adjustment of the free-convolution density matches the
normalizations of the two ensembles.
A uniform exponential bound for compact GOE counts and finite
inclusion-exclusion then upgrade correlation comparisons to the
Laplace functionals in Theorem~\ref{thm:main}.
\end{enumerate}

Two parts of this argument are useful beyond the final identification.
The entropy calculation gives Gaussian moment approximation with
polynomially small errors on common graph events of probability tending
to one, including graph-dependent choices of eigenvectors; its strictness
is a finite-dimensional tensor fact tied
to the open bulk and to \(d\ge3\).
The continuous switching argument converts a discrete finite-rank
symmetry into Gaussian perturbations of an exact limiting marked
process. The latter mechanism requires only fixed-dimensional
Gaussian information before the limit.

The passage from finitely many Gaussian marks to many weak switches
takes place after the graph-size limit. We first let \(n\to\infty\)
with fixed numbers of sampled edges and fixed spectral windows, then
control the tails to construct the principal-value marked law in
Section~\ref{sec:limits}. Its switching identities are exact for every
fixed number of channels. Only within this law do we let the channel
number \(k=2r\) tend to infinity, with the two spectral cutoffs prescribed
in Section~\ref{sec:switching}. Finally we increase the order of the
inclusion-exclusion brackets in Section~\ref{sec:dbm}. Thus the
finite-graph argument requires Gaussian approximation only in fixed
dimensions and at fixed moment orders.
